\section{Exact equality at the local CHSH bound}
\label{sec:ext-threshold}

The strict activation test has a complete companion classification at
equality. It is particularly useful when the generic matrix $K$ is
singular and the optimum is a rank-one limit.

\begin{theorem}[Three regimes at \(B_A^*=1\)]
\label{thm:ext-threshold}
Assume the Alice marginal is mixed and put $S=(B_A^*)^2$.
\begin{enumerate}[label=(\roman*)]
\item If $\alpha<1$, then $S\le1$, and
      $S=1$ if and only if $E=1$. Equality is never attained by an
      invertible filter.
\item If $\alpha=1$, then
      \[
       S=1\iff\det H_0=0
          \iff\det K=d^TKd.
      \]
      Equality is attained by an invertible filter if and only if
      $\ker H_0$ contains $z$ with $z^Tz\ne0$.
\item If $\alpha>1$, the optimum is attained by an invertible filter and
      \[
       S<1\iff R_A\succ0,\qquad
       S=1\iff R_A\succeq0,\ \det R_A=0,\qquad
       S>1\iff R_A\not\succeq0.
      \]
\end{enumerate}
The same statements hold for Bob after swapping the parties.
\end{theorem}

\begin{proof}
When $m<0$, the Gram identity and the compact reduction give
$f(n,v)\le1+m(1-\norm v^2)$.
Every interior pair has value strictly less than one. Hence a compact
maximizer of value one must lie on the boundary, where the maximum is
$E$. Equivalently, the actual-filter purity bound excludes finite
attainment because every invertible filter has positive determinant.

When $m=0$, $H_0=2G\succeq0$ and, using $\norm n=1$,
\begin{equation}
 f(n,v)=1-(n+iv)^\dagger H_0(n+iv).
 \label{eq:ext-threshold-form}
\end{equation}
If $H_0$ is nonsingular, its smallest eigenvalue is positive and the
subtracted term is uniformly positive on the compact domain.
If $0\ne z\in\ker H_0$, multiply $z$ by a phase so that
$\Re z\cdot\Im z=0$ and $\norm{\Re z}\ge\norm{\Im z}$.
The real part is then nonzero. Divide by its norm to get a feasible
pair $(n,v)$ with value one. The norm is strict exactly when the
original kernel contains a vector with $z^Tz\ne0$.
The identity $\det(K+iC)=\det K-d^TKd$ is polynomial in $K,d$:
it follows from diagonal $K$ and then orthogonal invariance, so it
also holds for singular $K$.

The final regime follows directly from
Theorem~\ref{thm:ext-shortcut}, since $\det K>0$ and
$R_A=\det K[M-(\tr A-1)I]$.
\end{proof}

\begin{proposition}[Boundary equality, product kernels, and rank]
\label{prop:ext-product-kernel}
For a two-qubit state with mixed filtered-side marginal,
\begin{equation}
 E=1\iff\ker\rho\ \text{contains a nonzero product vector}.
 \label{eq:ext-product-kernel}
\end{equation}
In particular, every state of rank at most two has $E=1$.
A rank-three state has $E=1$ precisely when its one-dimensional
kernel is generated by a product vector. At $\alpha=1$, a state of
rank at least three cannot attain $S=1$ by an invertible filter.
\end{proposition}

\begin{proof}
A rank-one Alice effect onto $|x\rangle$ produces the conditional
Bob matrix $\langle x|\rho|x\rangle$. It has positive trace because
$\rho_A$ is invertible. This $2\times2$ matrix is pure after
normalization exactly when it has a nonzero vector $|y\rangle$ in
its kernel. Positivity makes
$\langle x,y|\rho|x,y\rangle=0$ equivalent to
$|x\rangle\otimes|y\rangle\in\ker\rho$. The statement is unchanged
by invertible whitening, which maps product vectors to product vectors.

If $\rank\rho\le2$, the kernel contains a two-dimensional complex
subspace. Identify two independent vectors of this subspace with
$2\times2$ coefficient matrices $U,V$. The homogeneous quadratic
$\det(sU+tV)$ either vanishes identically or has a root in the complex
projective line. The associated nonzero matrix has rank one, and hence
represents a product vector. The rank-three assertion is immediate.

At $\alpha=1$, finite equality requires a nonisotropic vector in
$\ker H_0=\ker G$. The rank argument at the end of
Corollary~\ref{cor:ext-ceiling} excludes such a vector when the rank
is at least three.
\end{proof}

\begin{theorem}[Complete rational formula for rank at most two]
\label{cor:ext-low-rank}
If $\rank\rho\le2$ and the filtered-side marginal is mixed, then
\begin{equation}
 \boxed{\quad
 (B_A^*)^2=\max\left\{1,\frac{1+\alpha}{2}\right\}.
 \quad}
 \label{eq:ext-low-rank}
\end{equation}
In particular, the purity guard is both necessary and sufficient for
unilateral violation in this rank range.
The optimum is attained by an invertible filter exactly when
$\alpha\ge1$, and whitening alone attains it throughout that region.
For $\alpha>1$, whitening is the unique positive optimal filter up to
scale, with maximal success probability $1-\norm a$.

In the original state coordinates, let
$J_\rho=1+\norm{T}_F^2-\norm a^2-\norm b^2$ and
$G_A=\norm{T}_F^2+\norm a^2-\norm b^2-1$. Then
\begin{equation}
 \boxed{\quad
 (B_A^*)^2
 =\max\left\{1,\frac{J_\rho}{2(1-\norm a^2)}\right\}
 =1+\frac{\max(0,G_A)}{2(1-\norm a^2)}.
 \quad}
 \label{eq:ext-low-rank-raw}
\end{equation}
When both marginals are mixed, the larger local Bloch-vector length
gives a weakly larger unilateral optimum.
\end{theorem}

\begin{proof}
The canonical state is the normalized Choi matrix of a trace-preserving
qubit channel. Theorem 12 and equation (17) of
Ruskai--Szarek--Werner~\cite{ruskai2002} show that every such channel
of Choi rank at most two is unitarily equivalent to one with
translation $(0,0,\sin u\sin v)$ and diagonal linear part
$(\cos u,\cos v,\cos u\cos v)$.
This theorem includes the unital, non-extreme rank-two channels.
Accounting for the sign of $\sigma_y^T$ in the Choi correspondence,
and using local rotations to choose signs, our state therefore has
the form
\begin{equation}
 D=\diag(\sqrt x,-\sqrt y,\sqrt{xy}),\qquad
 c=(0,0,\sqrt{(1-x)(1-y)}),\qquad 0\le x,y\le1.
 \label{eq:ext-rank-two-normal-form}
\end{equation}
Thus $\alpha=2(x+y)-1$, and the canonical CHSH value is
$B(\rho_c)^2=x+y$, because $xy\le x,y$.

If $\alpha>1$, the purity ceiling is $(1+\alpha)/2=x+y$ and is
already achieved by the canonical state. If $\alpha=1$, the
purity bound is one and the canonical state again achieves it.
If $\alpha<1$, the purity obstruction gives the upper bound one,
and Proposition~\ref{prop:ext-product-kernel} supplies the matching
boundary value, with no finite attainment.
For $\alpha>1$, equality in~\eqref{eq:ext-filter-purity} forces
the optimal canonical effect to be $u=0$; hence the original
positive filter is proportional to $W$. Formula
\eqref{eq:ext-original-success} at $u=0$ gives $1-\norm a$.

Substitution of~\eqref{eq:ext-raw-alpha} proves
\eqref{eq:ext-low-rank-raw}. The common numerator satisfies
$J_\rho=4\tr(\rho\widetilde\rho)\ge0$.
Swapping the parties therefore changes only the positive denominator
$1-\norm a^2$ to $1-\norm b^2$, which proves the weak ordering.
If both ratios are at most one, the two optima are equal on the
plateau even when the local Bloch lengths differ.
\end{proof}

In particular, at $\alpha=1$ finite attainment of the local bound
is equivalent to $\rank\rho\le2$: the theorem proves sufficiency,
and Proposition~\ref{prop:ext-product-kernel} proves necessity.

For an explicit equality case of rank three, take
\begin{equation}
 \rho_{\mathrm{threshold}}=\frac18
 \begin{pmatrix}
 4&0&0&2\\0&0&0&0\\0&0&1&0\\2&0&0&3
 \end{pmatrix}.
 \label{eq:ext-threshold-example}
\end{equation}
Here $D=\diag(1/2,-1/2,3/4)$, $c=(0,0,1/4)$, and $\alpha=1$.
The eigenvalues of $H_0$ are $0,3/8,1/2$, and its kernel is spanned
by $(i,1,0)^T$, which is isotropic for $z^Tz$.
Thus $B_A^*=1$, but no invertible filter attains it.
The density-matrix kernel is generated by the product vector $|01\rangle$,
which also verifies $E=1$ directly.
